% !PRINTED 378
\weberpage{378}
%<<CONT
roots of (9) and only these, if in (8) one combines the roots
$x'$, $x''$ of (7) in pairs; and from this (6) follows.

By (6) it only remains to determine the function $\psi(D,m)$ for the
case that $m=q^k$ is a power of a prime $q$ that does not
divide $Q$. We need the symbol $(D,n)$ in
the sense explained in \S\,85, and we distinguish three cases:

1) If $q$ divides $\Delta$, then $(D,q)=0$, and the congruence
 (3) has a solution only when $k=1$; namely, according
as $D$ is even or odd, $x\equiv 0$ or $\equiv q \pmod{2q}$,
while there is no solution when $k>1$, and this holds also for
$q=2$. For in that case $Q$ would have to be odd and $\Delta$ divisible by 4,
consequently $x$ even. But the congruence

\[ \left(\frac{x}{2}\right)^2 \equiv Q^2\,\frac{\Delta}{4} \pmod{2^k} \]

is solvable only for $k=1$, since $\Delta:4$ is not a discriminant.

2) If $(D,q)=-1$, then the congruence (3) is not solvable,
because in that case $D$ is a quadratic non-residue of $q$ or (for $q=2$)
of 8.

3) If $(D,q)=+1$, then $D$ is a quadratic residue of $q$
(or of 8), and the congruence,

\[ x^2 \equiv D \mod 4q^k \]

has, as is known from number theory, \gesperrt{two roots}.

We collect these results in a survey as follows:

\[
\begin{array}{rcll}
\text{If }1) & (D,q)=0, & \text{then} & \psi(D,q)=1,\quad\psi(D,q^k)=0\quad(k>1),\\[0.4em]
\begin{array}{r}\text{\glqq\ }2)\\[0.3em]\text{\glqq\ }3)\end{array}
& \begin{array}{l}(D,q)=-1,\\[0.3em](D,q)=+1,\end{array}
& \begin{array}{c}\text{\glqq\ \glqq}\\[0.3em]\text{\glqq\ \glqq}\end{array}
& \left.\begin{array}{l}\psi(D,q^k)=0\\[0.3em]\psi(D,q^k)=2\end{array}\right\}\quad(k\geqq 1).
\end{array}
\]
% Source scan confirms the doubled greater-than-or-equal relation, transcribed as \geqq.

to which there is still to be added:

\[ \psi(D,1)=1. \]

By (5), therefore, $\psi(D,m)$ is always equal to zero when
a prime divides $m$ for which $(D,q)=-1$, or when
a prime factor of $D$ divides $m$ more than once.

In the other cases $\psi(D,q)$ is a power of 2 % [sic] German source prints q; context suggests m.
 whose
exponent is equal to the number of primes that divide $m$, but
 not $D$.

Between the symbols $\psi(D,m)$ and $(D,m)$ there subsists the
following general relation, always under the assumption
that $m$ is relatively prime to $Q$:
%CONT>>
